\documentclass[10pt,a4paper]{article}
\usepackage[T1]{fontenc}
\usepackage[utf8]{inputenc}
\usepackage{lmodern}
\usepackage{amsmath,amssymb}
\usepackage[margin=21mm]{geometry}
\usepackage{microtype}
\usepackage[hidelinks]{hyperref}
\usepackage{enumitem}
\setlist[itemize]{leftmargin=*,nosep,topsep=3pt}
\setlength{\parindent}{0pt}
\setlength{\parskip}{5pt}
\setlength{\emergencystretch}{3em}
\newcommand{\question}[1]{\par\medskip{\normalsize\bfseries #1}\par\smallskip}
\pagestyle{plain}

\begin{document}
{\Large\bfseries Photon Shot Noise Limits on Optical Detection of Neuronal Spikes and Estimation of Spike Timing\par}

\textbf{Paper:} Wilt, Fitzgerald and Schnitzer (2013),
\emph{Biophysical Journal} 104:51--62.\quad Theory, simulations and
comparison with published optical recordings.\par
\href{https://doi.org/10.1016/j.bpj.2012.07.058}{DOI: \nolinkurl{10.1016/j.bpj.2012.07.058}}\par
{\footnotesize Page references below use PDF page numbers.}

\question{To which general question does this paper relate?}
How accurately can optical recordings reveal neuronal spikes, and what
does this imply for estimating population firing-rate states? It concerns
measurement, not stimulation or synthesis of a neural controller.

\question{How does it answer the general questions?}
\textbf{Measured signal.} Calcium/voltage indicators produce fluorescence
from which spikes are inferred. Photon counts are Poisson-distributed:
their variance equals their mean. The single-spike model is
$F(t)=F_0+A e^{-t/\tau}\mathbf1_{t\ge0}$, where $F_0$ includes baseline
and unwanted background, $A$ is the fluorescence increment and $\tau$
its decay time (PDF pp.~1--2).

\textbf{Usable insight.} For small relative transients, sufficient photon
counts and sampling rate $\nu$ with $\nu\tau>1$, spike discriminability is
\[
 d'\approx\frac{\Delta F}{F}\sqrt{\frac{F_0\tau}{2}}
       =A\sqrt{\frac{\tau}{2F_0}}.
\]
Larger relative responses and less unwanted background improve detection.
Longer transients collect more photons but impair timing and separation
of nearby spikes. Faster sampling alone does not ensure better estimates.
Likelihood-based detection quantifies missed-spike/false-alarm tradeoffs
(PDF pp.~3--7, 9--11).

\textbf{Timing and feedback.} Detection and timing are distinct tasks.
For $d'=5$, $\tau=0.15$~s and $\nu=20$~Hz, the calculated timing-error
standard-deviation lower bound is 20.7~ms. This bound applies to unbiased
estimators and need not be attained; it is not an online processing-delay
guarantee (PDF pp.~6--7, 10).

\textbf{Model relevance (derived).} A possible population-rate observation
approximation, with $r$ the firing-rate state and $c$ an indicator state, is
\[
 \dot c=-c/\tau+q r,\qquad
 Y_k\sim\operatorname{Poisson}\!\left(
   \int_{t_k}^{t_{k+1}}[F_0+\alpha c(t)]\,dt\right).
\]
This motivates an observer that includes indicator dynamics and counting
noise; it does not establish observability. Averaging reconstructed spikes
from representative neurons can estimate a population rate. The paper
does not validate this population model or direct full-state feedback.

\textbf{Limits.} The main analysis concerns isolated/sparse spikes,
linear unsaturated responses, known waveforms and already separated cells.
Motion, optical mixing, saturation and overlapping spikes can worsen
performance. The bounds are ideal limits; technology comparisons describe
2013 devices, rather than current optical capability (PDF pp.~2, 7, 9--11).

\question{What papers does it cite that could provide more information?}
\begin{itemize}
\item Yaksi and Friedrich (2006), \emph{Reconstruction of firing rate changes
across neuronal populations by temporally deconvolved Ca$^{2+}$ imaging}
(ref.~24): directly relevant to population-rate reconstruction.
\item Vogelstein et al. (2010), \emph{Fast nonnegative deconvolution for
spike train inference from population calcium imaging} (ref.~18).
\item Grewe et al. (2010), \emph{High-speed in vivo calcium imaging reveals
neuronal network activity with near-millisecond precision} (ref.~15).
\end{itemize}

\end{document}
